\begin{statementbox}
	\textbf{\textcolor{CStatement}{条件}}
	\begin{description}[style=nextline, leftmargin=2.8em, labelsep=0.5em, itemsep=0.35em]
		\item[\textbf{\textcolor{CStatement}{条件1（凸性条件）：}}] 函数$f$在区间$I$上严格凸（或严格凹）。
		\item[\textbf{\textcolor{CStatement}{条件2（变量条件）：}}] $x_i\in I$，权重$\lambda_i>0$且$\sum_{i=1}^n\lambda_i=1$。
	\end{description}

	\textbf{\textcolor{CStatement}{结论}}
	\begin{description}[style=nextline, leftmargin=2.8em, labelsep=0.5em, itemsep=0.45em]
		\item[\textbf{\textcolor{CStatement}{结论一（琴生不等式）：}}]
			\textbf{\textcolor{CStatement}{直观描述：}}
			凸函数加权平均的函数值不大于函数值的加权平均。
			\[
				f\!\left(\sum_{i=1}^n\lambda_i x_i\right)\le\sum_{i=1}^n\lambda_i f(x_i)
			\]
			等号成立当且仅当$x_1=x_2=\cdots=x_n$.

		\item[\textbf{\textcolor{CStatement}{推广结论（严格凹情形）：}}]
			\textbf{\textcolor{CStatement}{直观描述：}}
			严格凹时，加权平均的函数值不小于函数值的加权平均。
			\[
				f\!\left(\sum_{i=1}^n\lambda_i x_i\right)\ge\sum_{i=1}^n\lambda_i f(x_i)
			\]
			等号成立当且仅当$x_1=x_2=\cdots=x_n$.
	\end{description}
\end{statementbox}
